\subsection{NODAL GROUP OF A TORUS}

The nodal group of a torus S, denoted by \(\Gamma(\mathrm{S})\) , is the kernel of the exponential map \(\mathrm{L}(\mathrm{S}) \to \mathrm{S}\) . This is a discrete subgroup of \(\mathrm{L}(\mathrm{S})\) , whose rank is equal to the dimension of S, and the R-linear map \(\mathbf{R} \otimes_{\mathbf{Z}} \Gamma(\mathrm{S}) \to \mathrm{L}(\mathrm{S})\) that extends the canonical injection of \(\Gamma(\mathrm{S})\) into \(\mathrm{L}(\mathrm{S})\) is bijective. It induces by passage to the quotient an isomorphism \(\mathbf{R}/\mathbf{Z} \otimes_{\mathbf{Z}} \Gamma(\mathrm{S}) \to \mathrm{S}\) .

For example, the nodal group \(\Gamma(\mathbf{U})\) of U is the subgroup \(2\pi iZ\) of \(\mathrm{L}(\mathbf{U}) = i\mathbf{R}\) .

For any morphism of tori \(f: \mathrm{S} \to \mathrm{S}'\) , denote by \(\Gamma(f)\) the homomorphism \(\Gamma(\mathrm{S}) \to \Gamma(\mathrm{S}')\) induced by \(\mathrm{L}(f)\) . We have a commutative diagram

\[
\begin{array}{c c c c c c c c c} 0 & \longrightarrow & \Gamma (\mathrm{S}) & \longrightarrow & \mathrm{L} (\mathrm{S}) & \stackrel {{\exp_ {\mathrm{S}}}} {{\longrightarrow}} & \mathrm{S} & \longrightarrow & 0 \\ & & \Big \downarrow^ {\Gamma (f)} & & \Big \downarrow^ {\mathrm{L} (f)} & & \Big \downarrow^ {f} \\ 0 & \longrightarrow & \Gamma (\mathrm{S} ^ {\prime}) & \longrightarrow & \mathrm{L} (\mathrm{S} ^ {\prime}) & \stackrel {{\exp_ {\mathrm{S} ^ {\prime}}}} {{\longrightarrow}} & \mathrm{S} ^ {\prime} & \longrightarrow & 0. \end{array}\tag{1}
\]

Let \(a \in \mathrm{X}(\mathrm{S})\) ; applying the preceding to the morphism from S to U defined by \(a\) , we see that the C-linear map \(\delta(a) : \mathrm{L}(\mathrm{S})_{(\mathbf{C})} \to \mathbf{C}\) of no. 1 maps \(\Gamma(\mathrm{S})\) to \(2\pi i\mathbf{Z}\) . Thus, we can define a Z-bilinear form on \(\mathrm{X}(\mathrm{S}) \times \Gamma(\mathrm{S})\) by putting

\[
\langle a, X \rangle = \frac {1}{2 \pi i} \delta (a) (X), \quad a \in \mathrm{X(S)}, X \in \Gamma (\mathrm{S}).\tag{2}
\]

PROPOSITION 3. The bilinear form \((a,X)\mapsto \langle a,X\rangle\) on \(\mathrm{X(S)}\times \Gamma (\mathrm{S})\) is invertible.

Recall (Algebra, Chap. IX) that, by definition, this means that the linear maps \(\mathrm{X}(\mathrm{S}) \to \mathrm{Hom}_{\mathbf{Z}}(\varGamma(\mathrm{S}), \mathbf{Z})\) and \(\varGamma(\mathrm{S}) \to \mathrm{Hom}_{\mathbf{Z}}(\mathrm{X}(\mathrm{S}), \mathbf{Z})\) associated to this bilinear form are bijective.

It is immediate that if the conclusion of the proposition is true for two tori, it is also true for their product. Thus, since every torus of dimension n is isomorphic to \(U^{n}\) , we are reduced to the case in which S = U. In this particular case, the assertion is immediate.

Let \(f: S \to S'\) be a morphism of tori. Then, each of the linear maps \(X(f): X(S') \to X(S)\) and \(\Gamma(f): \Gamma(S) \to \Gamma(S')\) is the transpose of the other: for all \(a' \in X(S')\) and all \(X \in \Gamma(S)\) ,

\[
\langle X (f) (a ^ {\prime}), X \rangle = \langle a ^ {\prime}, \Gamma (f) (X) \rangle .\tag{3}
\]

PROPOSITION 4. Let S and S' be tori. Denote by M(S,S') the group of morphisms of Lie groups from S to S'. The maps \(f \mapsto \mathrm{X}(f)\) and \(f \mapsto \Gamma(f)\) are isomorphisms of groups from M(S,S') to Hom \(_{\mathbf{Z}}(\mathrm{X}(\mathrm{S}'),\mathrm{X}(\mathrm{S}))\) and to Hom \(_{\mathbf{Z}}(\Gamma(\mathrm{S}),\Gamma(\mathrm{S}'))\) , respectively.

If \(f\) is a morphism of Lie groups from S to \(S'\) , the homomorphism \(X(f)\) is simply the dual of \(f\) in the sense of Spectral Theories, Chap. II, §1, no. 7. The map \(\varphi \mapsto \hat{\varphi}\) from \(\mathrm{Hom}_{\mathbf{Z}}(\mathrm{X}(\mathrm{S}'), \mathrm{X}(\mathrm{S}))\) to \(\mathrm{M}(\mathrm{S},\mathrm{S}')\) defined in loc. cit. is the inverse of the map \(f \mapsto \mathrm{X}(f)\) from \(\mathrm{M}(\mathrm{S},\mathrm{S}')\) to \(\mathrm{Hom}_{\mathbf{Z}}(\mathrm{X}(\mathrm{S}'), \mathrm{X}(\mathrm{S}))\); the latter is thus bijective. If we identify \(\Gamma(\mathrm{S})\) (resp. \(\Gamma(\mathrm{S}')\) ) with the dual \(\mathbf{Z}\) -module of \(\mathrm{X}(\mathrm{S})\) (resp. \(\mathrm{X}(\mathrm{S}')\) ) (Prop. 3), \(\Gamma(f)\) coincides with the transpose of the homomorphism \(X(f)\) , hence the proposition.

Remarks. 1) Let \(f: \mathrm{S} \to \mathrm{S}'\) be a morphism of tori. The snake diagram (Algebra, Chap. X, §1, no. 2) associated to (1) gives an exact sequence

\[
0 \longrightarrow \mathrm{Ker} \Gamma (f) \longrightarrow \mathrm{KerL} (f) \longrightarrow \mathrm{Ker} f \stackrel {d} {\longrightarrow}\tag{4}
\]

\[
\stackrel {{d}} {{\longrightarrow}} \operatorname{Coker} \Gamma (f) \longrightarrow \operatorname{Coker} \mathrm{L} (f) \longrightarrow \operatorname{Coker} f \longrightarrow 0.
\]

In particular, assume that \(f\) is surjective, with finite kernel \(\mathbf{N}\) , so that we have an exact sequence

\[
0 \longrightarrow \mathrm{N} \stackrel {i} {\longrightarrow} \mathrm{S} \stackrel {f} {\longrightarrow} \mathrm{S} ^ {\prime} \longrightarrow 0,
\]

where i is the canonical injection. Then, \(L(f)\) is bijective, and (4) gives an isomorphism \(N \to \text{Coker } \Gamma(f)\) , hence an exact sequence

\[
0 \longrightarrow \Gamma (\mathrm{S}) \stackrel {\Gamma (f)} {\longrightarrow} \Gamma (\mathrm{S} ^ {\prime}) \longrightarrow \mathrm{N} \longrightarrow 0.\tag{5}
\]

Moreover, by Spectral Theories, Chap. II, §1, no. 7, Th. 4, the sequence

\[
0 \longrightarrow \mathrm{X} (\mathrm{S} ^ {\prime}) \stackrel {{\mathrm{X} (f)}} {{\longrightarrow}} \mathrm{X} (\mathrm{S}) \stackrel {{\mathrm{X} (i)}} {{\longrightarrow}} \mathrm{X} (\mathrm{N}) \longrightarrow 0\tag{6}
\]

is exact.

\begin{enumerate}
\def\labelenumi{\arabic{enumi})}
\setcounter{enumi}{1}
\tightlist
\item
  By Prop. 4, the map \(f \mapsto \Gamma(f)(2\pi i)\) from \(\mathrm{M}(\mathbf{U},\mathrm{S})\) to \(\Gamma(\mathrm{S})\) is an isomorphism; if \(a \in \mathrm{X}(\mathrm{S}) = \mathrm{M}(\mathrm{S},\mathbf{U})\) and \(f \in \mathrm{M}(\mathbf{U},\mathrm{S})\) , then the composite \(a \circ f \in \mathrm{M}(\mathbf{U},\mathbf{U})\) is the endomorphism \(u \mapsto u^r\) , where \(r = \langle a,\Gamma(f)(2\pi i)\rangle\) . We shall identify \(\mathrm{M}(\mathbf{U},\mathbf{U}) = \mathrm{X}(\mathbf{U})\) with \(\mathbf{Z}\) from now on, the element \(r\) of \(\mathbf{Z}\) being associated to the endomorphism \(u \mapsto u^r\) ; thus, with the notations above,
\end{enumerate}

\[
a \circ f = \langle a, \Gamma (f) (2 \pi i) \rangle .
\]

\begin{enumerate}
\def\labelenumi{\arabic{enumi})}
\setcounter{enumi}{2}
\item
  To the exact sequence \(0 \to \Gamma(\mathrm{S}) \to \mathrm{L}(\mathrm{S}) \xrightarrow{\exp_{\mathrm{S}}} \mathrm{S} \to 0\) is associated an isomorphism from \(\Gamma(\mathrm{S})\) to the fundamental group of \(\mathrm{S}\) , called canonical in the sequel. For any morphism of tori \(f: \mathrm{S} \to \mathrm{S}'\) , \(\Gamma(f)\) can then be identified via the canonical isomorphisms \(\Gamma(\mathrm{S}) \to \pi_1(\mathrm{S})\) and \(\Gamma(\mathrm{S}') \to \pi_1(\mathrm{S}')\) with the homomorphism \(\pi_1(f): \pi_1(\mathrm{S}) \to \pi_1(\mathrm{S}')\) induced by \(f\) . Note that this gives another interpretation of the exact sequence (5) (cf.~General Topology, Chap. XI, in preparation).
\item
  The homomorphisms of \(\mathbf{Z}\) -modules \(\delta : \mathrm{X}(\mathrm{S}) \to \mathrm{Hom}_{\mathbf{C}}(\mathrm{L}(\mathrm{S})_{(\mathbf{C})}, \mathbf{C})\) and \(\iota : \Gamma(\mathrm{S}) \to \mathrm{L}(\mathrm{S})_{(\mathbf{C})}\) ( \(\iota\) is induced by the canonical injection of \(\Gamma(\mathrm{S})\) into \(\mathrm{L}(\mathrm{S})\) ) extend to isomorphisms of \(\mathbf{C}\) -vector spaces
\end{enumerate}

\[
\begin{array}{l} u: \mathbf {C} \otimes \mathrm{X} (\mathrm{S}) \to \mathrm{Hom} _ {\mathbf {C}} (\mathrm{L} (\mathrm{S}) _ {(\mathbf {C})}, \mathbf {C}) \\ v: \mathbf {C} \otimes \Gamma (\mathrm{S}) \to \mathrm{L} (\mathrm{S}) _ {(\mathbf {C})} \end{array}
\]

which we shall call canonical in the sequel. Note that, if we extend the pairing between \(X(S)\) and \(\Gamma(S)\) by C-linearity to a bilinear form \(\ll\) , \(\gg\) on \((\mathbf{C} \otimes \mathrm{X}(S)) \times (\mathbf{C} \otimes \Gamma(S))\) , then

\[
\langle u (a), v (b) \rangle = 2 \pi i \ll a, b \gg .
\]

